% \iffalse meta-comment
%
% File: pkuthss3.dtx Copyright (C) 2025 Xuchen Lin
%
% This file is part of `pkuthss3', a LaTeX3-restructured version of
% `pkuthss' by Casper Ti. Vector et al.
% -------------------------------------------
%
% It may be distributed and/or modified under the
% conditions of the LaTeX Project Public License, either version 1.3c
% of this license or (at your option) any later version.
% The latest version of this license is in
%    https://www.latex-project.org/lppl.txt
% and version 1.3c or later is part of all distributions of LaTeX
% version 2008 or later.
%
% This file has the LPPL maintenance status "author-maintained".
%
% The files belonging to `pkuthss3' are pkuthss3.ins and pkuthss3.dtx,
% and the derived files pkuthss3.cls and pkuthss3.pdf.
%
%
% \fi
%
% \iffalse
%<cls>\NeedsTeXFormat{LaTeX2e}
%<*driver>
\ProvidesFile{pkuthss3.dtx}[%
%</driver>
%<cls>\ProvidesExplClass {pkuthss3}
  {2025/04/18} {0.0}
%<cls>  {Peking University dissertation document class}
%<*driver>
]
\documentclass{l3doc}

\usepackage[UTF8,linespread=1.1,fontset=ubuntu]{ctex}
\newcommand\pkuthss{PKU\-\textsc{thss}3}

\EnableCrossrefs
\CodelineIndex
\setlength\IndexMin{200pt}

\GetFileInfo{\jobname.dtx}
\begin{document}
  \DocInput{\jobname.dtx}
\end{document}
%</driver>
% \fi
%
% \title{\pkuthss}
% \author{
%   Xuchen Lin
% }
% \date{Version \fileversion, Released \filedate}
%
% \maketitle
%
% \tableofcontents
%
% \begin{documentation}
%
% \end{documentation}
%
% \begin{implementation}
% \clearpage
% \section{代码实现 Implementation}
%
%    开始文档类内容。定义内部宏前缀为 \cls{pkuthss}。
%    \begin{macrocode}
%<*cls>
%<@@=pkuthss>
%    \end{macrocode}
%
%    定义 \cls{pkuthss} 的信息类别为文档类。
%    \begin{macrocode}
\prop_gput:Nnn \g_msg_module_type_prop { pkuthss } { Class }
%    \end{macrocode}
%
% \subsection{\LaTeX{} 内核版本检查 \LaTeX{} version check}
%    本文档类使用了一些 2025/06/01 版本内核的功能，因此需要判断内核版本是否够新。
%
% \begin{macro}[int]{\IfFormatAtLeastTF}
%    内核版本判定。本宏在 2020/10/01 版本加入，因此若缺少，则提供定义。
%    \begin{macrocode}
\providecommand \IfFormatAtLeastTF { \@ifl@t@r \fmtversion }
%    \end{macrocode}
% \end{macro}
%
%    \begin{macrocode}
\IfFormatAtLeastTF { 2024-06-01 }
  {}
  {
    \ClassError { pkuthss3 } { LaTeX~kernel~too~old }
      {
        You~need~a~newer~LaTeX~to~use~this~release~of~pkuthss3.\\
        使用本版本~pkuthss3~需要更新版本的~LaTeX。
      }
    \endinput
  }
%    \end{macrocode}
%
% \subsection{内部函数与变量 Internal functions and variables}
%
% \begin{macro}{\l_@@_word_grid_size_dim}
%    选项用到的变量。
%    \begin{macrocode}
\dim_new:N \l_@@_word_grid_size_dim
%    \end{macrocode}
% \end{macro}
%
% \begin{macro}
%   {
%     \@@_define_option:n, \@@_define_metadata:n,
%     \@@_set_option:n,    \@@_set_metadata:n
%   }
%    用于便捷定义或设置文档类选项（option）与元数据（metadata）的简称。
%    \begin{macrocode}
\cs_new_protected:Nn \@@_define_option:n
  { \keys_define:nn { pkuthss }            {#1} }
\cs_new_protected:Nn \@@_define_metadata:n
  { \keys_define:nn { pkuthss / metadata } {#1} }
\cs_new_protected:Nn \@@_set_option:n
  { \keys_set:nn    { pkuthss }            {#1} }
\cs_new_protected:Nn \@@_set_metadata:n
  { \keys_set:nn    { pkuthss / metadata } {#1} }
%    \end{macrocode}
% \end{macro}
%
% \subsubsection{错误与警告 Errors and warnings}
%    \begin{macrocode}
\msg_new:nnn { pkuthss } { option-invalid }
  { Option~`\l_keys_key_tl'~is~invalid. }
\msg_new:nnn { pkuthss } { option-preset }
  {
    Option~`\l_keys_key_tl'~has~been~set~by~pkuthss.\\
    User-provided~value~is~ignored.
  }
%    \end{macrocode}
%
% \subsection{选项定义 Defining options}
%    \begin{macrocode}
\@@_define_option:n
  {
%    \end{macrocode}
%
% \subsubsection{文档类加载时选项 Class loading-time options}
% \begin{macro}
%   {
%     language, \l_@@_main_language_tl,
%     traditional-chinese, \l_@@_traditional_chinese_bool,
%     english-variant, \l_@@_english_variant_tl
%   }
%    设置论文正文的\textbf{主要}语言，会据此导入 \pkg{babel} 并影响“目录”“参考文献”等章节
%    标题的本地化，影响西文单词的断字，还会影响无文档网格时“单倍行距”按照中易宋体还是 Times
%    New Roman 设置。
%
%    北大似乎只允许中英日韩德俄法西八种语言的学位论文，因此只允许选择这些语言。中日韩语会按照
%    中易宋体设置无网格单倍行距，其他语言按照 Times New Roman。
%
%    如果需要章节序号“第一章”等不按照正文语言本地化，请更改 chapter-localization 选项。
%
%    默认值为 chinese。中文的简体繁体需另行通过 traditional-chinese 选项设置。
%
%    英文变体通过 english-variant 指定。
%
%    这里将输入选项保存为 token list，因为众多 \pkg{babel} 的实现使用了 \cs{ifx} 进行
%    语言判断，而存为 str 会把 catagory code 全部转成 12，导致 \cs{ifx} 失败。
%    \begin{macrocode}
    language .choices:nn =
      {
        chinese, english, japanese, korean,
        german,  ngerman, russian,  french, spanish
      }
      {
        \tl_set_eq:NN \l_@@_main_language_tl \l_keys_choice_tl
      },
    language .initial:n = chinese,
    language .usage:n   = load,
    traditional-chinese .bool_set:N = \l_@@_traditional_chinese_bool,
    traditional-chinese .initial:n  = false,
    traditional-chinese .default:n  = true,
    traditional-chinese .usage:n    = load,
    english-variant .choices:nn =
      { american, australian, british, canadian, newzealand }
      {
        \tl_set_eq:NN \l_@@_english_variant_tl \l_keys_choice_tl
      },
    english-variant .initial:n  = american,
    english-variant .usage:n    = load,
%    \end{macrocode}
% \end{macro}
%
% \subsubsection{文档类一般选项 Class general options}
% \begin{macro}{thesis-type, \l_@@_thesis_type_str}
%    论文类型。
%    \begin{macrocode}
    thesis-type .choices:nn =
      { bachelor, master, doctor, postdoc }
      {
        \str_set:NV \l_@@_thesis_type_str \l_keys_choice_tl
      },
    thesis-type .initial:n = doctor,
%    \end{macrocode}
% \end{macro}
%
% \begin{macro}{word-grid, \l_@@_word_grid_bool, word-grid-size}
%    是否启用类似 MS Word 文档网格的行为，以及若启用，以 bp 为单位给出网格间距的值。
%    \begin{macrocode}
    word-grid      .bool_set:N = \l_@@_word_grid_bool,
    word-grid      .initial:n  = true,
    word-grid-size .code:n     =
      {
        \fp_set:Nn \l_tmpa_fp {#1}
        \dim_set:Nn \l_@@_word_grid_size_dim
          { \fp_eval:n \l_tmpa_fp bp }
      },
    word-grid-size .initial:n = 15.6,
%    \end{macrocode}
% \end{macro}
%
% \subsubsection{其他未知选项 Other unknown options}
%    使用 \cs[no-index]{IfClassLoadedTF} 判定该未知选项是否是在加载本文档类时给出的。若
%    为加载时选项，则传给上游 \cls{ctexbook}，否则给出警告。
%    \begin{macrocode}
    unknown .code:n =
      {
        \IfClassLoadedTF { ctexbook }
          { \msg_warning:nn { pkuthss } { option-invalid } }
          { \PassOptionsToClass { \CurrentOption } { ctexbook } }
      },
%    \end{macrocode}
%
% \subsubsection{^^A
%   从标准文档类或 C\TeX{} 中去除的选项 Options removed from standard or C\TeX\
%   classes}
%    以下两个 C\TeX{} 选项 \pkg{pkuthss} 已经给定值，用户另行指定值无效，因此给出警告。
%    \begin{macrocode}
    zihao      .code:n = \msg_warning:nn { pkuthss } { option-preset },
    linespread .code:n = \msg_warning:nn { pkuthss } { option-preset },
%    \end{macrocode}
%
%    \begin{macrocode}
  }
%    \end{macrocode}
%
%    以下标准文档类选项由于格式要求，对 \pkg{pkuthss} 无意义，设置为无效。
%    \begin{macrocode}
\clist_map_inline:nn
  {
    a5paper , b5paper , letterpaper , legalpaper , executivepaper ,
    landscape ,
    10pt , 11pt , 12pt ,
    twoside , oneside ,
    titlepage , notitlepage ,
    onecolumn , twocolumn ,
    leqno ,
    openbib ,
%    \end{macrocode}
%    以及 C\TeX{} 的 \cmd[no-index]{scheme} 选项。
%    \begin{macrocode}
    scheme
  }
  {
    \@@_define_option:n
      {
        #1 .code:n  = \OptionNotUsed,
        #1 .usage:n = load,
      }
  }
%    \end{macrocode}
%
%
% \subsection{加载上游文档类与宏包 Loading upstream class and packages}
%    处理选项后加载 \cls{ctexbook}。
%    \begin{macrocode}
\ProcessKeyOptions [ pkuthss ]
\LoadClass [ a4paper , zihao = -4 , linespread = 1.0 ] { ctexbook }
%    \end{macrocode}
%
% \subsubsection{\pkg{babel} 语言设置 \pkg{babel} settings}
%
%    总是需要导入英文。
%    \begin{macrocode}
\PassOptionsToPackage { \l_@@_english_variant_tl } { babel }
%    \end{macrocode}
%    主要语言不是英文时，还需导入主要语言。
%    \begin{macrocode}
\str_if_eq:VnF \l_@@_main_language_tl { english }
  {
    \PassOptionsToPackage { \l_@@_main_language_tl } { babel }
  }
%    \end{macrocode}
%    主要语言是中文或朝鲜/韩文时，需要将主要语言以“modern”形式导入。
%    \begin{macrocode}
\str_case:VnT \l_@@_main_language_tl
  {
    { chinese } { }
    { korean }  { }
  }
  { \PassOptionsToPackage { provide = * } { babel } }
%    \end{macrocode}
%    导入 \pkg{babel}。
%    \begin{macrocode}
\RequirePackage { babel }
%    \end{macrocode}
%    主要语言非中文或英文时，需要将中文作为次要语言导入。
%    \begin{macrocode}
\str_case:VnF \l_@@_main_language_tl
  {
    { chinese } { }
    { english } { }
  }
  { \AtBeginDocument { \babelprovide [ import ] { chinese } } }
%    \end{macrocode}
%
% \subsection{用户设置接口及元数据 User settings API and metadata}
% \begin{macro}{\pkusetup,\pkumetadata}
%    简单设置为别名。设置为只可在导言区使用。
%    \begin{macrocode}
\NewDocumentCommand \pkusetup { m }
  { \@@_set_option:n   {#1} }
\NewDocumentCommand \pkumetadata { m }
  { \@@_set_metadata:n {#1} }
\@onlypreamble \pkusetup
\@onlypreamble \pkumetadata
%    \end{macrocode}
% \end{macro}
%
%
% \subsection{本地化 Localization}
%    大多数本地化通过 \pkg{babel} 自动完成。
%
% \subsubsection{繁体中文 Traditional Chinese}
%    \begin{macrocode}
\bool_if:NT \l_@@_traditional_chinese_bool
  {
    \setlocalecaption { chinese } { ref }        { 參考文獻 }
    \setlocalecaption { chinese } { bib }        { 參考文獻 }
    \setlocalecaption { chinese } { appendix }   { 附錄 }
    \setlocalecaption { chinese } { contents }   { 目錄 }
    \setlocalecaption { chinese } { listfigure } { 插圖 }
    \setlocalecaption { chinese } { figure }     { 圖 }
    \setlocalecaption { chinese } { proof }      { 證明 }
  }
%    \end{macrocode}
%
% \subsubsection{朝鲜/韩文 Korean}
%    需要保留谚文之间的空格。
%    \begin{macrocode}
\str_if_eq:VnF \l__pkuthss_main_language_tl { korean }
  { \AtBeginDocument { \babelprovide [ import ] { korean } } }
\AtBeginDocument
  {
    \addto \extraskorean   { \ctexset { space = true } }
    \addto \noextraskorean { \ctexset { space = auto } }
  }
%    \end{macrocode}
%
%    \begin{macrocode}
%</cls>
%    \end{macrocode}
%
% \end{implementation}
%
% \PrintIndex
% \PrintChanges
